% Part: proof-theory
% Chapter: normalization
% Section: reductions

\documentclass[../../../include/open-logic-section]{subfiles}

\begin{document}

\olfileid[hi]{pt}{nor}{red}

\olsection{अपचयन रूपांतरण}

यदि \Log{N1i} के किसी !!{proof} में कट की लंबाई~$1$ है, तो वह किसी एक
!!{formula}-आवृत्ति वाला खंड है, और वह !!{formula} किसी !!a{introduction}
अनुमान का निष्कर्ष तथा किसी !!a{elimination} अनुमान का मुख्य आधार-वाक्य दोनों
है। प्रसामान्यीकरण प्रमाण में ऐसे कटों का ``अपचयन'' उस
!!{introduction}/!!{elimination} युग्म पर समाप्त होने वाले उप-!!{proof} को ऐसे
नए !!{proof} से बदलकर किया जाता है जिसमें अनुमानों का वह युग्म हटा दिया गया
हो। ऐसे प्रतिस्थापन से नए कट बन सकते हैं। किंतु हम देख सकते हैं कि नए कटों की
कोटि अपचयित किए जा रहे कट से कम होती है। अतः नए !!{proof} की या तो कट लंबाई
कम है (क्योंकि लंबाई~$1$ वाला कोई महत्तम कट हट गया), या उसकी कट कोटि कम हो गई
है (यदि संबंधित महत्तम कट ही एकमात्र महत्तम कट था)।

उदाहरणतः यदि संबंधित अनुमान \Intro{\lif} और \Elim{\lif} हैं, $!A \ident !B
\lif !C$, और कोई उप-!!{proof}~$\delta_1$ निम्न पर समाप्त होता है:
\[
\AxiomC{$\Discharge{!B}{x}$}
\RightLabel{$\delta_2$}
\DeduceC{$!C$}
\DischargeRule{\Intro{\lif}}{x}
\UnaryInfC{$!B \lif !C$}
\AxiomC{}
\RightLabel{$\delta_3$}
\DeduceC{$!B$}
\RightLabel{\Elim{\lif}}
\BinaryInfC{$!C$}
\DisplayProof
\]
इस कट को हटाने के लिए हम~$\delta$ में उपप्रमाण~$\delta_1$ को निम्न से बदलते
हैं:
\[
\AxiomC{}
\RightLabel{$\delta_3$}
\DeduceC{$!B$}
\RightLabel{$\Subst{\delta_2}{\delta_3}{!B^x}$}
\DeduceC{$!C$}
\DisplayProof
\]
$\Subst{\delta_2}{\delta_3}{!B^x}$ से हमारा आशय उस !!{proof} से है जो
~$\delta_2$ में~$x$ से चिह्नित प्रत्येक मान्यता~$!B$ को पूरे
!!{proof}~$\delta_3$ से बदलने पर प्राप्त होता है। इस !!{proof} में अब~$x$ से
चिह्नित~$!B$ रूप की खुली मान्यताएँ नहीं हैं। इसमें~$x$ से चिह्नित कोई दूसरी
मान्यता भी नहीं है, क्योंकि सदा की तरह हम मानते हैं कि कोई दो !!{formula}s
एक ही चिह्न वाली मान्यताओं के रूप में नहीं आते। अतः नए !!{proof} की खुली
मान्यताएँ~$\delta_1$ जैसी ही हैं;~$\delta_2$ का मान्यताएँ मुक्त करने वाला हर
अनुमान $\Subst{\delta_2}{\delta_3}{!B^x}$ में उन्हीं मान्यताओं को मुक्त करता
है; और~$\delta_3$ का मान्यताएँ मुक्त करने वाला हर अनुमान संभवतः अनेक प्रतियों
में परिणत होता है, जो $\Subst{\delta_2}{\delta_3}{!B^x}$ में संगत मान्यताओं
को मुक्त करती हैं।

$\delta_2$, $\delta_3$ या $\delta_4$ के पूर्णतः भीतर, अथवा~$\delta$ में
~$\delta_1$ के पूर्णतः बाहर स्थित कोई भी कट अपरिवर्तित रहता है: उसमें अब भी
वही !!{formula}s और वही अनुमान सम्मिलित हैं, इसलिए विशेषतः उसकी कोटि और लंबाई
~$\delta$ के संगत कट जैसी ही हैं। हमने लंबाई~$1$ का कोई महत्तम कट हटाया है;
इसलिए या तो कट लंबाई कम हुई है, अथवा कट कोटि कम हुई है (यदि~$\delta$ की कट
लंबाई~$1$ थी और यही एकमात्र महत्तम कट था)।

किंतु दो बातें हो सकती हैं:
\begin{enumerate}
  \item सभी मान्यताओं $!B^x$ को $\delta_3$ से बदलने पर हमने~$\delta_3$ के
  सभी कटों की प्रतियाँ बढ़ा दी हैं। यदि उनमें से कुछ महत्तम कट थे, तो नए
  !!{proof} की कट लंबाई पुराने !!{proof} की कट लंबाई से अधिक हो सकती है। हमें
  इसकी रोकथाम सुनिश्चित करनी है। इसके लिए अपचयन रूपांतरण केवल \emph{सबसे
  दाएँ} महत्तम कटों पर लगाते हैं; अर्थात् ऐसे कटों पर जिनसे दाईं ओर~$\delta$
  से गुजरती किसी शाखा पर कोई महत्तम कट न हो। यदि~$\delta_3$ में कोई कट होता,
  तो वह यहाँ अपचयित किए जा रहे कट के दाईं ओर होता और हम सबसे दाएँ कट पर काम
  नहीं कर रहे होते। सदा सबसे दाएँ कट चुनकर हम सुनिश्चित करते हैं कि या तो
  !!{proof} की कट लंबाई घटेगी, या कट कोटि भी घट जाएगी।
  \item यदि $\delta_2$ की कोई मान्यता $!B^x$ किसी !!a{elimination} अनुमान का
  मुख्य आधार-वाक्य है, और~$\delta_3$ का अंतिम अनुमान \Elim{\lor},
  \Elim{\lexists}, या कोई !!a{introduction} अनुमान है, तो ऐसी मान्यता पर
  ~ $\delta_3$ को $\delta_2$ में कलम करके हमने नया कट बना दिया है। जो
  ~ $\delta_2$ में केवल मान्यता थी वह अब लंबाई~$1$ का कट
  (!!a{introduction} अनुमान का निष्कर्ष और !!a{elimination} अनुमान का मुख्य
  आधार-वाक्य), अथवा किसी कट खंड की अंतिम !!{formula}-आवृत्ति
  (!!a{elimination} अनुमान का मुख्य आधार-वाक्य और~\Elim{\lor} या
  \Elim{\lexists} का निष्कर्ष) है। यदि $B^x$, \Elim{\lor} या
  \Elim{\lexists} का गौण आधार-वाक्य था, तो वह पहले से किसी खंड, और संभवतः कट
  खंड, की पहली !!{formula}-आवृत्ति था। नए !!{proof} में वह खंड अब अधिक लंबा
  हो सकता है। किंतु ध्यान दें कि इस नए कट या विद्यमान कट खंड को बनाने वाला
  !!{formula}~$!B$ है और $\depth{!B} < \depth{!B \lif !C}$। अतः यद्यपि हमने
  संभवतः कट जोड़े या उनकी लंबाई बढ़ाई है, इनमें से कोई महत्तम कट नहीं है।
  इसलिए या तो कट कोटि घटी है, या समान कोटि के कट शेष हों तो कट लंबाई घटी है।
\end{enumerate}

लंबाई~$1$ के कट के अपचयन की प्रक्रिया को \emph{अपचयन रूपांतरण} (या
\emph{चक्कर रूपांतरण}) कहते हैं। कुछ नियम-युग्मों के प्रतिरूप
\olref{tab:reduction-conversions} में दिए गए हैं। (शेष अभ्यास के लिए छोड़े गए
हैं।)

\begin{table}
\begin{multline*}
\bottomAlignProof
\AxiomC{$\Discharge{!B}{x}$}
\RightLabel{$\delta_2$}
\shortDeduce
\DeduceC{$!C$}
\DischargeRule{\Intro{\lif}}{x}
\UnaryInfC{$!B \lif !C$}
\AxiomC{}
\RightLabel{$\delta_3$}
\shortDeduce
\DeduceC{$!B$}
\RightLabel{\Elim{\lif}}
\BinaryInfC{$!C$}
\DisplayProof
\quad \longrightarrow\quad
\shoveright{\bottomAlignProof
\AxiomC{}
\RightLabel{$\delta_3$}
\shortDeduce
\DeduceC{$!B$}
\RightLabel{$\Subst{\delta_2}{\delta_3}{!B^x}$}
\shortDeduce
\DeduceC{$!C$}
\DisplayProof}
\\[4ex]
\shoveleft{\bottomAlignProof
\AxiomC{}
\RightLabel{$\delta_2$}
\shortDeduce
\DeduceC{$!B$}
\AxiomC{}
\RightLabel{$\delta_3$}
\shortDeduce
\DeduceC{$!C$}
\RightLabel{\Intro{\land}}
\BinaryInfC{$!B \land !C$}
\RightLabel{\Elim{\land}[1]}
\UnaryInfC{$!B$}
\DisplayProof
\quad \longrightarrow\quad
\bottomAlignProof
\AxiomC{}
\RightLabel{$\delta_2$}
\shortDeduce
\DeduceC{$!B$}
\DisplayProof}\qquad
\shoveright{\bottomAlignProof
\AxiomC{}
\RightLabel{$\delta_2$}
\shortDeduce
\DeduceC{$!B$}
\AxiomC{}
\RightLabel{$\delta_3$}
\shortDeduce
\DeduceC{$!C$}
\RightLabel{\Intro{\land}}
\BinaryInfC{$!B \land !C$}
\RightLabel{\Elim{\land}[2]}
\UnaryInfC{$!C$}
\DisplayProof
\quad \longrightarrow\quad
\bottomAlignProof
\AxiomC{}
\RightLabel{$\delta_3$}
\shortDeduce
\DeduceC{$!C$}
\DisplayProof}
\\[4ex]
\bottomAlignProof
\AxiomC{}
\RightLabel{$\delta_2$}
\shortDeduce
\DeduceC{$!B$}
\RightLabel{\Intro{\lor}[1]}
\UnaryInfC{$!B \lor !C$}
\AxiomC{$\Discharge{!B}{x}$}
\RightLabel{$\delta_3$}
\shortDeduce
\DeduceC{$!D$}
\AxiomC{$\Discharge{!C}{y}$}
\RightLabel{$\delta_4$}
\shortDeduce
\DeduceC{$!D$}
\DischargeRule{\Elim{\lor}}{x\,y}
\TrinaryInfC{$!D$}
\DisplayProof
\quad \longrightarrow\quad
\shoveright{\bottomAlignProof
\AxiomC{}
\RightLabel{$\delta_2$}
\shortDeduce
\DeduceC{$!B$}
\RightLabel{$\Subst{\delta_3}{\delta_2}{!B^x}$}
\shortDeduce
\DeduceC{$!D$}
\DisplayProof}
\\[2ex]
\shoveleft{\bottomAlignProof
\AxiomC{}
\RightLabel{$\delta_2$}
\shortDeduce
\DeduceC{$!B(c)$}
\RightLabel{\Intro{\lforall}}
\UnaryInfC{$\lforall[x][!B(x)]$}
\RightLabel{\Elim{\lforall}}
\UnaryInfC{$!B(t)$}
\DisplayProof
\quad \longrightarrow\quad
\bottomAlignProof
\AxiomC{}
\RightLabel{$\Subst{\delta_2}{t}{c}$}
\shortDeduce
\DeduceC{$!B(t)$}
\DisplayProof}
\\[4ex]
\shoveright{\bottomAlignProof
\AxiomC{}
\RightLabel{$\delta_2$}
\shortDeduce
\DeduceC{$!B(t)$}
\RightLabel{\Intro{\lexists}}
\UnaryInfC{$\lexists[x][!B(x)]$}
\AxiomC{$\Discharge{!B(c)}{x}$}
\RightLabel{$\delta_3$}
\shortDeduce
\DeduceC{$!D$}
\DischargeRule{\Elim{\lexists}}{x}
\BinaryInfC{$!D$}
\DisplayProof
\quad \longrightarrow\quad
\bottomAlignProof
\AxiomC{}
\RightLabel{$\delta_2$}
\shortDeduce
\DeduceC{$!B(t)$}
\RightLabel{$\Subst{\Subst{\delta_3}{t}{c}}{\delta_2}{!B(t)^x}$}
\shortDeduce
\DeduceC{$!D$}
\DisplayProof}
\end{multline*}
\caption{\Log{N1i} के कुछ अपचयन रूपांतरण।}
\ollabel{tab:reduction-conversions}
\end{table}

\begin{prob}
\Intro{\lnot} के बाद आने वाले \Elim{\lnot} के लिए अपचयन रूपांतरण दीजिए।
\end{prob}

एक स्थिति पर अभी विचार शेष है। कट की परिभाषा में वह स्थिति भी है जहाँ कट की
लंबाई~$1$ है और $!A$,~\FalseInt{} का निष्कर्ष तथा किसी !!a{elimination}
अनुमान का मुख्य आधार-वाक्य दोनों है। इस स्थिति से उसी प्रकार निपटते हैं जैसे
उस स्थिति से जिसमें $!A$ किसी !!a{introduction} नियम का निष्कर्ष हो। उदाहरणतः
निम्न को बदलें
\[
\bottomAlignProof
\AxiomC{}
\RightLabel{$\delta_2$}
\DeduceC{$\lfalse$}
\UnaryInfC{$!B \lif !C$}
\AxiomC{}
\RightLabel{$\delta_3$}
\DeduceC{$!B$}
\RightLabel{\Elim{\lif}}
\BinaryInfC{$!C$}
\DisplayProof
\quad\text{से}\quad
\bottomAlignProof
\AxiomC{}
\RightLabel{$\delta_2$}
\DeduceC{$\lfalse$}
\UnaryInfC{$!C$}
\DisplayProof
\]
उन स्थितियों में हमें थोड़ा सावधान रहना होगा जहाँ $!A \ident (!B \lor !C)$
या $!A \ident \lexists[x][!B(x)]$। उदाहरणतः यदि हमने निम्न को
\[
\bottomAlignProof
\AxiomC{}
\RightLabel{$\delta_2$}
\DeduceC{$\lfalse$}
\RightLabel{\FalseInt}
\UnaryInfC{$!B \lor !C$}
\AxiomC{$\Discharge{!B}{x}$}
\RightLabel{$\delta_3$}
\DeduceC{$!D$}
\AxiomC{$\Discharge{!C}{x}$}
\RightLabel{$\delta_4$}
\DeduceC{$!D$}
\RightLabel{\Elim{\lor}}
\TrinaryInfC{$!D$}
\DisplayProof
\quad\text{से}\quad
\bottomAlignProof
\AxiomC{}
\RightLabel{$\delta_2$}
\DeduceC{$\lfalse$}
\RightLabel{\FalseInt}
\UnaryInfC{$!D$}
\DisplayProof
\]
बदला, तो हम कट !!{formula}~$!D$ वाला नया कट प्रस्तुत कर सकते हैं, और इसकी
कोई गारंटी नहीं कि $\depth{!D} < \depth{!B \lor !C}$!{} इसलिए यहाँ भी हमें
उसी प्रकार आगे बढ़ना चाहिए जैसे \Intro{\lor}/\Elim{\lor} अपचयन में बढ़े थे,
और निम्न से बदलना चाहिए
\[
\bottomAlignProof
\AxiomC{}
\RightLabel{$\delta_2$}
\DeduceC{$\lfalse$}
\RightLabel{\FalseInt}
\UnaryInfC{$!B$}
\RightLabel{$\delta_3$}
\DeduceC{$!D$}
\DisplayProof
\quad\text{अथवा}\quad
\bottomAlignProof
\AxiomC{}
\RightLabel{$\delta_2$}
\DeduceC{$\lfalse$}
\RightLabel{\FalseInt}
\UnaryInfC{$!C$}
\RightLabel{$\delta_4$}
\DeduceC{$!D$}
\DisplayProof
\]

\end{document}
